\section*{Références bibliographiques et crédits}
\textit{Théorie des intersections et théorème de Riemann--Roch},
Séminaire de Géométrie Algébrique du Bois-Marie 1966--67 (SGA~6), séminaire
dirigé par P.~Berthelot, A.~Grothendieck et L.~Illusie, avec la collaboration
de D.~Ferrand, J.-P.~Jouanolou, O.~Jussila, S.~Kleiman, M.~Raynaud et
J.-P.~Serre. Lecture Notes in Mathematics~225, Springer, 1971.

\subsection*{Crédits de la collection conservés d'après les documents fournis}
Les volumes originaux de SGA conservent leurs auteurs, les participants au
séminaire et leurs éditeurs, nommés dans les textes correspondants. La
collection comprend des travaux d'Alexander Grothendieck, Michèle Raynaud,
Michel Raynaud, Pierre Deligne, Nicholas Katz, Michel Demazure, Michael Artin,
Jean-Louis Verdier, Luc Illusie et de leurs collaborateurs; cette liste ne
remplace pas les crédits détaillés de chaque volume.

SGA~1 conserve la nouvelle composition française associée au projet bénévole
de Bas Edixhoven, à la Société mathématique de France et aux notes de Michel
Raynaud, d'après arXiv math/0206203v2.

SGA~2 est AI-translated à partir du TeX français corrigé de la réédition
dirigée par Yves Laszlo, \textit{Cohomologie locale des faisceaux cohérents et
théorèmes de Lefschetz locaux et globaux (SGA~2)}, Documents Mathématiques~4,
Société mathématique de France, 2005 (arXiv math/0511279v1), dont la
numérotation et les étiquettes sont conservées. La majeure partie de la
composition française initiale est due aux bénévoles du projet de Bas
Edixhoven : L.~Bayle, N.~Borne, O.~Brinon, J.~Buresi, M.~Chardin, F.~Ducrot,
P.~Graftiaux, F.~Han, P.~Karwasz, L.~Koelblen, D.~Madore, S.~Morel, D.~Naie,
B.~Osserman, J.~Riou et V.~Sécherre. C.~Sabbah l'a adaptée au format de la SMF.

Les sources anglaises de SGA~2 et SGA~3 créditent la traduction anglaise de
Jacob C.~Reinhold, version \texttt{e7a259f3f8608ad3edf9bf6eead3fd504dd2d23e},
utilisée pour la comparaison et la rédaction. La mention CC BY~4.0 qu'il a
indiquée s'applique à cette contribution et est conservée; elle ne concède
pas de licence sur les œuvres françaises sous-jacentes ni sur toutes les
contributions de cette collection. Les formulations anglaises héritées
restent attribuées à leur source humaine. Dans SGA~3, les lemmes,
démonstrations et corrections ajoutés dans la réédition de Philippe
Gille--Patrick Polo sont présentés dans les notes du traducteur et les
additions qui la citent.

SGA~4 est AI-translated à partir de l'édition française en LaTeX de SGA~4,
tomes~1--3, distribuée par Fabrice Orgogozo (version de la source au commit
git \texttt{71766d9}), et utilise ses fichiers de configuration TeX. Les notes
de bas de page intitulées « N.D.E. » sont les notes éditoriales de cette
édition, traduites.

SGA~4\textonehalf{} conserve la composition LaTeX de Daniel Miller et sa
note originale de composition. Cette composition reste créditée à Daniel
Miller. Les documents fournis indiquent que son dépôt
\url{https://github.com/dkmiller/sga4.5} n'est plus en ligne, que le texte
français figure dans sga-fr et qu'une copie subsiste dans NomiL/sga4.5.

Les traductions anglaises supplémentaires sont AI-translated; les corrections,
la révision, la reconstruction et la composition ultérieures sont produites
par l'IA. L'éditeur n'a pas traduit le texte. Les crédits humains propres aux
sources et les véritables mentions de droits des tiers restent applicables;
l'attribution n'établit ni permission ni approbation.

